% Part: turing-machines
% Chapter: machines-computations
% Section: variants

\documentclass[../../../include/open-logic-section]{subfiles}

\begin{document}

\olfileid{tur}{mac}{var}
\olsection{ట్యూరింగ్ యంత్రాల రూపాంతరాలు}

ట్యూరింగ్ యంత్రాలను నిర్వచించడానికి నిజానికి అనేక సాధ్యమైన పద్ధతులు ఉన్నాయి;
మనది వాటిలో ఒకటి మాత్రమే. కొన్ని విషయాల్లో మన నిర్వచనం ఇతర నిర్వచనాల కంటే
స్వేచ్ఛాయుతమైనది. మనం ఏ పరిమిత వర్ణమాలనైనా అనుమతిస్తాం; మరింత నియంత్రితమైన
నిర్వచనం $\TMstroke$, $\TMblank$ అనే రెండు టేపు సంకేతాలను మాత్రమే అనుమతించవచ్చు.
యంత్రం టేపుపై ఒక సంకేతాన్ని వ్రాస్తూనే కదలడాన్ని మనం అనుమతిస్తాం; ఇతర
నిర్వచనాలు వ్రాయడం లేదా కదలడం---వీటిలో ఒకదాన్నే అనుమతిస్తాయి. టేపు శీర్షాన్ని
కదపకుండా వ్రాసే అవకాశాన్ని మనం అనుమతిస్తాం; ఇతర నిర్వచనాలు $\TMstay$
``నిర్దేశాన్ని'' విడిచిపెడతాయి. మరికొన్ని విషయాల్లో మన నిర్వచనం మరింత
నియంత్రితమైనది. టేపు ఒకే దిశలో అనంతంగా ఉంటుందని ఊహించాం; ఇతర నిర్వచనాలు
టేపు ఎడమకు, కుడికి రెండు వైపులా అనంతంగా ఉండటాన్ని అనుమతిస్తాయి. నిజానికి,
ఎన్ని వేర్వేరు టేపులనైనా, లేదా గడుల అనంత జాలకాన్నైనా అనుమతించవచ్చు. ట్యూరింగ్
యంత్రపు నిర్దేశ సమితికి మనం సంక్రమణ ప్రమేయంతో ప్రాతినిధ్యం ఇస్తాం; ఇతర
నిర్వచనాలు సంక్రమణ సంబంధాన్ని ఉపయోగిస్తాయి, అందులో యంత్రానికి ఏ ఇచ్చిన
పరిస్థితిలోనైనా ఒకటి కంటే ఎక్కువ సాధ్యమైన నిర్దేశాలు ఉంటాయి.

నిర్వచనానికి చేసిన ఈ చివరి సడలింపు ప్రత్యేకంగా ఆసక్తికరమైనది. మన నిర్వచనంలో,
యంత్రం స్థితి~$q$లో సంకేతం~$\sigma$ను చదువుతున్నప్పుడు, కొత్త సంకేతం, స్థితి,
టేపు శీర్షం స్థానం ఏమిటో $\delta(q, \sigma)$ నిర్ణయిస్తుంది. కానీ నిర్దేశ
సమితి ప్రస్తుత స్థితి--సంకేత జతలు $\tuple{q, \sigma}$కు, కొత్త
స్థితి--సంకేతం--దిశ త్రయాలు $\tuple{q', \sigma', D}$కు మధ్య సంబంధంగా ఉండటాన్ని
అనుమతిస్తే, ట్యూరింగ్ యంత్రపు చర్య ఏకైకంగా నిర్ణయించబడకపోవచ్చు---నిర్దేశ
సంబంధంలో $\tuple{q, \sigma, q', \sigma', D}$,
$\tuple{q, \sigma, q'', \sigma'', D'}$ రెండూ ఉండవచ్చు. ఈ సందర్భంలో మనకు
\emph{అనిర్ణీత} ట్యూరింగ్ యంత్రం ఉంటుంది. గణనా సంక్లిష్టతా సిద్ధాంతంలో ఇవి
ముఖ్యమైన పాత్ర పోషిస్తాయి.

ట్యూరింగ్ యంత్రం ఎప్పుడు ఆగుతుందనే విషయంలో కూడా వేర్వేరు సంప్రదాయాలు ఉన్నాయి:
సంక్రమణ ప్రమేయం నిర్వచితం కానప్పుడు అది ఆగుతుందని మనం అంటాం; ఇతర నిర్వచనాలు
యంత్రం ప్రత్యేకంగా నిర్దేశించిన ఆగే స్థితిలో ఉండాలని కోరతాయి. నిర్దేశిత ఆగే
స్థితిని కోరడం ట్యూరింగ్ యంత్రాలు గణించగలవాటిని ప్రభావితం చేసే నియంత్రణ కాదని
\olref[hal]{sec}లో వివరించాం. మన ట్యూరింగ్ యంత్రాల టేపులు ఒకే దిశలో అనంతంగా
ఉంటాయి కనుక, ట్యూరింగ్ యంత్రం ఒక నిర్దేశాన్ని సరిగా అమలు చేయలేని సందర్భాలు
ఉంటాయి: అది అత్యంత ఎడమ గడిని చదువుతూ ఎడమకు కదలవలసి ఉంటే. మన నిర్వచనం
ప్రకారం, అది ``కింద పడిపోకుండా'' ఉన్నచోటే ఉంటుంది; కానీ అలా జరిగినప్పుడు
ఆగిపోయేలా కూడా నిర్వచించి ఉండవచ్చు. ఈ నిర్వచనం కూడా సమానమే: గడి~$0$ నుండి
ఎడమకు కదలడానికి ప్రయత్నించినప్పుడు ఆగే ట్యూరింగ్ యంత్రపు ప్రవర్తనను అనుకరించడానికి,
ప్రతి సంక్రమణ $\delta(q,\TMendtape) =
\tuple{q',\sigma,\TMleft}$ను తొలగించవచ్చు---అప్పుడు $\TMendtape$పై ఎడమకు
కదలడానికి ప్రయత్నించడానికి బదులుగా యంత్రం ఆగుతుంది.\footnote{ఇది
\emph{పూర్తిగా} పనిచేయదు; ఎందుకంటే గడి~$0$ కాక ఇతర గడులపై $\TMendtape$ను
వ్రాయడం, చదవడం నుండి మనల్ని ఏదీ నిరోధించదు (\olref[una]{ex:mover} చూడండి).
అలాంటి ప్రయోజనానికి బదులుగా ఉపయోగించడానికి రెండవ~$\TMendtape'$ సంకేతాన్ని
చేర్చి ఈ సమస్యను పరిష్కరించవచ్చు.}

సంఖ్యలకు ప్రాతినిధ్యం వహించే పద్ధతులు (తద్వారా ట్యూరింగ్ యంత్రం గణించే
ఇన్‌పుట్--అవుట్‌పుట్ ప్రమేయాలు) కూడా వేర్వేరుగా ఉంటాయి: మనం ఏకాంక
ప్రాతినిధ్యాన్ని ఉపయోగిస్తాం; ద్వ్యాంక ప్రాతినిధ్యాన్ని కూడా ఉపయోగించవచ్చు.
దీనికి $\TMblank$, $\TMendtape$లకు అదనంగా రెండు సంకేతాలు కావాలి.

ఇప్పుడు ఒక ఆసక్తికరమైన వాస్తవం: ఏ ప్రమేయాలు ట్యూరింగ్-గణనీయాలు అనే విషయంలో
ఈ రూపాంతరాల్లో ఏదీ తేడా కలిగించదు. \emph{ఒక నిర్వచనం ప్రకారం ఒక ప్రమేయం
ట్యూరింగ్-గణనీయమైతే, అది వాటన్నిటి ప్రకారమూ ట్యూరింగ్-గణనీయమే.}

దీనిని నిర్ధారించే వివరాల్లోకి మనం వెళ్లం. ఒకే ఉదాహరణ ఇక్కడ ఉంది: రెండు
దిశల్లో అనంతమైన టేపును లేదా అనేక టేపులను అనుమతించడం వల్ల మనకు అదనపు గణన
సామర్థ్యం రాదు. ఎందుకంటే, ఒకే దిశలో అనంతమైన ఒక్క టేపు గల ట్యూరింగ్ యంత్రం
అనేక టేపులను లేదా రెండు దిశల్లో అనంతమైన టేపులను అనుకరించగలదు. ఉదాహరణకు,
అదనపు స్థితులు, నిర్దేశాలను ఉపయోగించి, అనేక టేపులు లేదా రెండు దిశల్లో అనంతమైన
టేపు గల యంత్రపు ప్రోగ్రామును, ఒకే దిశలో అనంతమైన ఒక్క టేపు గల యంత్రపు
ప్రోగ్రాముగా ``అనువదించవచ్చు.'' అనువదించిన యంత్రం సరి గడులను టేపు~$1$ యొక్క
గడులకు (లేదా రెండు దిశల్లో అనంతమైన టేపు యొక్క ``ధన'' గడులకు), బేసి గడులను
టేపు~$2$ యొక్క గడులకు (లేదా ``ఋణ'' గడులకు) ఉపయోగించవచ్చు.

\end{document}
